\documentclass[11pt,a4paper]{report}

%==========================================================
% Packages
%==========================================================
\usepackage[utf8]{inputenc}
\usepackage[T1]{fontenc}
\usepackage{lmodern}  % uncomment for nicer fonts if available
\usepackage[margin=1in]{geometry}
\usepackage{amsmath,amssymb,amsthm,amsfonts}
\usepackage{mathtools}
\usepackage{enumitem}
\usepackage{hyperref}
\hypersetup{colorlinks=true, linkcolor=black, urlcolor=blue, citecolor=black}
\usepackage{microtype}  % uncomment if you have lmodern or another scalable font

\setlength{\parindent}{0pt}
\setlength{\parskip}{0.6em}

\theoremstyle{plain}
\newtheorem{theorem}{Theorem}[chapter]
\newtheorem{lemma}[theorem]{Lemma}
\newtheorem{proposition}[theorem]{Proposition}
\newtheorem{corollary}[theorem]{Corollary}

\theoremstyle{definition}
\newtheorem{definition}[theorem]{Definition}
\newtheorem{example}[theorem]{Example}
\newtheorem{remark}[theorem]{Remark}
\newtheorem{algorithm}[theorem]{Algorithm}

\newcommand{\todo}{\textit{[Proof to be filled in.]}}

\newcommand{\R}{\mathbb{R}}
\newcommand{\N}{\mathbb{N}}
\newcommand{\norm}[1]{\left\lVert#1\right\rVert}
\newcommand{\inner}[2]{\langle #1,\,#2\rangle}
\newcommand{\grad}{\nabla}
\newcommand{\hess}{\nabla^{2}}
\DeclareMathOperator*{\argmin}{arg\,min}
\DeclareMathOperator*{\argmax}{arg\,max}
\DeclareMathOperator{\diag}{diag}
\DeclareMathOperator{\rank}{rank}
\newcommand{\trans}{^{\top}}
\newcommand{\eps}{\varepsilon}

\title{C6.2 \\[0.3em] \Large Continuous Optimization}
\author{Lecture notes \\ \small (based on the lectures of Coralia Cartis, Oxford)}
\date{}

\begin{document}

\maketitle

\noindent\textbf{References.}
\begin{itemize}[leftmargin=2em,itemsep=0pt]
  \item N.~I.~M.~Gould, \emph{An Introduction to Algorithms for Continuous Optimization},
        2006.
  \item J.~Nocedal and S.~J.~Wright, \emph{Numerical Optimization}, Springer (2nd ed.~2006).
  \item A.~R.~Conn, N.~I.~M.~Gould and Ph.~L.~Toint, \emph{Trust-Region Methods}, SIAM 2000.
  \item R.~Fletcher, \emph{Practical Methods of Optimization}, 2nd ed., Wiley 1987.
\end{itemize}

\tableofcontents

%==========================================================
\chapter{Optimality Conditions for Unconstrained Optimization}
%==========================================================

\section{Problem classes and notation}

We consider the unconstrained optimization problem
\begin{equation}\tag{UP}\label{eq:UP}
\min_{x\in\R^{n}}\; f(x),
\end{equation}
where $f:\R^{n}\to\R$ is smooth (typically $f\in\mathcal{C}^{1}$ or $\mathcal{C}^{2}$).

\paragraph{Iteration notation.} Iterates and other vector sequences are indexed by
superscripts, e.g.\ $x^{k}$, $s^{k}$, $u^{k}$, $\delta^{k}$, so that subscripts
remain available for components such as $x_i$ and $\lambda_i$. Scalar algorithmic
parameters and named model quantities are indexed by subscripts, e.g.\ $\alpha_k$,
$\eps_k$, $\Delta_k$, $\rho_k$, $B_k$, $m_k$.

\begin{definition}[Local and global minimizers]
A point $x^{*}\in\R^{n}$ is a \emph{local minimizer} of $f$ if there exists
$\delta>0$ such that $f(x^{*})\le f(x)$ for all $x\in N(x^{*},\delta)
=\{x:\norm{x-x^{*}}<\delta\}$. It is a \emph{strict} local minimizer if the
inequality is strict for $x\ne x^{*}$, and a \emph{global minimizer} if
$f(x^{*})\le f(x)$ for all $x\in\R^{n}$.
\end{definition}

\begin{definition}[Descent direction]
Given $x\in\R^{n}$ with $\grad f(x)\ne 0$, a vector $s\in\R^{n}$ is a
\emph{descent direction} for $f$ at $x$ if $\grad f(x)\trans s<0$.
\end{definition}

\begin{lemma}[Descent property]\label{lem:descent}
If $f\in\mathcal{C}^{1}$ and $s$ is a descent direction at $x$, then there
exists $\bar{\alpha}>0$ such that $f(x+\alpha s)<f(x)$ for all
$\alpha\in(0,\bar{\alpha}]$.
\end{lemma}
\begin{proof}
  Since $s$ is a descent direction, $\grad f(x)\trans s<0$. By continuity of
  $\grad f$, there exists $\bar{\alpha}>0$ such that
  \[
    \grad f(x+t s)\trans s<0
  \]
  for all $t\in[0,\bar{\alpha}]$. Hence, for any $\alpha \in (0, \bar{\alpha}]$, by the mean value theorem there exists $t\in(0, \alpha)$ such that
  \[
    f(x+\alpha s)-f(x) = \alpha \cdot \grad f(x+t s)\trans s < 0.
  \]
\end{proof}

\section{First-order optimality conditions}

\begin{theorem}[First-order necessary conditions]\label{thm:foc}
Let $f\in\mathcal{C}^{1}(\R^{n})$. If $x^{*}$ is a local minimizer of $f$,
then $\grad f(x^{*})=0$.
\end{theorem}
\begin{proof}
  Suppose not. Put $s = -\grad f(x^{*})$ and note that $s$ is a descent direction. By Lemma~\ref{lem:descent} there exists $\bar{\alpha}>0$ such that $f(x^{*}+\alpha s)<f(x^{*})$ for all $\alpha \in (0, \bar{\alpha}]$, contradiction.
\end{proof}

\begin{definition}[Stationary point]
A point $x$ with $\grad f(x)=0$ is called a \emph{stationary point} of $f$.
\end{definition}

\section{Second-order optimality conditions}

\begin{theorem}[Second-order necessary conditions]\label{thm:soc-nec}
Let $f\in\mathcal{C}^{2}(\R^{n})$. If $x^{*}$ is a local minimizer of $f$,
then $\grad f(x^{*})=0$ and $\hess f(x^{*})\succeq 0$ (positive semidefinite),
i.e.\ $s\trans\hess f(x^{*})s\ge 0$ for all $s\in\R^{n}$.
\end{theorem}
\begin{proof}
  By Theorem~\ref{thm:foc}, $\grad f(x^{*})=0$. Suppose for contradiction that $s\trans\hess f(x^{*})s < 0$ for some $s \in \R^{n}$. Since $s \neq 0$ and $f \in \mathcal{C}^2$, there exists $\hat{\alpha} > 0$ such that $s\trans\hess f(x^{*} + \alpha s)s < 0$ for all $\alpha \in (0, \hat{\alpha}]$. By Taylor's theorem, for any $\alpha \in (0, \hat{\alpha}]$, there exists $\bar{\alpha} \in (0, \alpha)$ such that
  \[
    f(x^{*} + \alpha s) - f(x^{*}) = \frac{\alpha^2}{2} s\trans\hess f(x^{*} + \bar{\alpha} s)s < 0,
  \]
  contradiction.
\end{proof}

\begin{theorem}[Second-order sufficient conditions]\label{thm:soc-suf}
Let $f\in\mathcal{C}^{2}(\R^{n})$. If $\grad f(x^{*})=0$ and $\hess f(x^{*})\succ 0$
(positive definite), then $x^{*}$ is a strict local minimizer of $f$.
\end{theorem}
\begin{proof}
  Since $f \in \mathcal{C}^2$, there exists $\hat{\alpha} > 0$ such that $s\trans\hess f(x^{*} + \alpha s)s > 0$ for all $\alpha \in (0, \hat{\alpha}]$ and $s \in \R^{n}$. By Taylor's theorem, for any $\alpha \in (0, \hat{\alpha}]$, there exists $\bar{\alpha} \in (0, \alpha)$ such that
  \[
    f(x^{*} + \alpha s) - f(x^{*}) = \frac{\alpha^2}{2} s\trans\hess f(x^{*} + \bar{\alpha} s)s > 0.
  \]
  Thus $x^{*}$ is a strict local minimizer within the ball $N(x^{*}, \hat{\alpha})$.
\end{proof}

\begin{example}[Necessary $\not\Rightarrow$ sufficient]
The conditions are different:
\begin{itemize}[leftmargin=2em,itemsep=0pt]
  \item $f(x)=x^{3}$, $x^{*}=0$: satisfies the second-order necessary condition
        ($f''(0)=0\ge 0$) but is not a local minimizer.
  \item $f(x)=x^{4}$, $x^{*}=0$: is a strict local minimizer but the second-order
        sufficient condition \emph{fails} ($f''(0)=0\not>0$).
\end{itemize}
The gap between necessary and sufficient is the case $\hess f(x^{*})\succeq 0$ but
not $\succ 0$, in which case optimality cannot be decided from $\grad f$ and
$\hess f$ alone.
\end{example}

\section{Convex problems}

For convex problems, the picture simplifies dramatically: \emph{every local
minimizer is global}, and the first-order necessary condition $\grad f(x^{*})=0$
becomes \emph{sufficient}. This is one of the main reasons convex optimization is
algorithmically tractable.

\begin{definition}[Convex function]
A function $f:\R^{n}\to\R$ is \emph{convex} if its domain is convex and, for all
$x,y$ in its domain and all $\lambda\in[0,1]$,
$f(\lambda x+(1-\lambda)y)\le \lambda f(x)+(1-\lambda)f(y)$.
\end{definition}

\begin{theorem}[Optimality for convex problems]\label{thm:convex-opt}
Let $f:\R^{n}\to\R$ be convex and $\mathcal{C}^{1}$. Then $x^{*}$ is a global
minimizer of $f$ if and only if $\grad f(x^{*})=0$.
\end{theorem}
\begin{proof}
  If $x^{*}$ is a global minimizer, then by Theorem~\ref{thm:foc}, $\grad f(x^{*})=0$. It remains to show the converse.

  We first show that if $x^{*}$ is a local minimizer, then it is a global minimizer. Suppose not. There exists $y \in \R^{n}$ such that $f(y) < f(x^{*})$. Put $s = y - x^{*}$. By convexity, for $\lambda \in (0, 1)$,
  \[
    f(x^{*} + \lambda s) = f((1 - \lambda)x^* + \lambda y) \leq (1 - \lambda) f(x^{*}) + \lambda f(y) < f(x^{*}).
  \]
  But then $x^*$ is a local minimizer, so there exists $\delta > 0$ such that $f(x^{*} + \lambda s) \leq f(x^{*})$ for all $\lambda \in (0, \delta)$, contradiction.

  \todo
\end{proof}


%==========================================================
\chapter{Linesearch Methods}
%==========================================================

A linesearch method generates iterates $x^{k+1}=x^{k}+\alpha_{k}s^{k}$ where
$s^{k}$ is a descent direction and $\alpha_{k}>0$ is a stepsize. The philosophy is
``liberal in the choice of search direction, keeping bad behaviour in control by
choice of $\alpha_{k}$''. Choosing $\alpha_{k}$ properly is essential: too long and
the method may overshoot, too short and progress is negligible.

\section{The generic linesearch method}

\begin{algorithm}[Generic Linesearch Method (GLM)]\label{alg:GLM}
Choose $\eps>0$ and $x^{0}\in\R^{n}$. For $k\ge 0$:
\begin{enumerate}[leftmargin=2em,itemsep=0pt]
  \item If $\norm{\grad f(x^{k})}\le\eps$, stop.
  \item Compute a descent direction $s^{k}$ with $\grad f(x^{k})\trans s^{k}<0$.
  \item Compute a stepsize $\alpha_{k}>0$ along $s^{k}$ such that
        $f(x^{k}+\alpha_{k}s^{k})<f(x^{k})$.
  \item Set $x^{k+1}:=x^{k}+\alpha_{k}s^{k}$.
\end{enumerate}
\end{algorithm}

\section{Stepsize rules: the Armijo condition}

\paragraph{Choices of $\alpha_{k}$.} The conceptually simplest choice is the
\emph{exact linesearch}
$\alpha_{k}:=\argmin_{\alpha>0} f(x^{k}+\alpha s^{k})$, but this is computationally
expensive for nonlinear $f$. Inexact linesearches replace it with a cheap rule that
guarantees the stepsize is ``not too long'' and ``not too short''.

The \emph{Armijo condition} demands sufficient decrease: $\alpha$ is acceptable
provided the actual decrease is at least a fraction $\beta$ of what the linear model
predicts. In practice $\beta$ is taken small, e.g.\ $\beta=0.1$ or even $\beta=0.001$,
to keep the rule mild.

For $\beta\in(0,1)$ the \emph{Armijo (sufficient decrease) condition} at
iteration $k$ is
\begin{equation}\label{eq:armijo}\tag{ac}
f(x^{k}+\alpha s^{k})\le f(x^{k})+\beta\alpha\,\grad f(x^{k})\trans s^{k}.
\end{equation}

\begin{remark}[Wolfe conditions]
The Armijo condition prevents $\alpha_{k}$ from being too long, but on its own
admits arbitrarily small $\alpha_{k}$. The \emph{Wolfe conditions} add a second
(``curvature'') inequality
$\grad f(x^{k}+\alpha s^{k})\trans s^{k}\ge \eta\,\grad f(x^{k})\trans s^{k}$
with $\eta\in(\beta,1)$, ruling out steps that are too short. Wolfe linesearches
are needed for some methods (e.g.\ BFGS, see Chapter~\ref{ch:qn-nls}) to ensure
the curvature condition $(\delta^{k})\trans\gamma^{k}>0$.
\end{remark}

\begin{algorithm}[Backtracking-Armijo (bArmijo) linesearch]\label{alg:bArmijo}
Choose $\alpha^{(0)}>0$, $\tau\in(0,1)$, $\beta\in(0,1)$. Set $i:=0$.
While \eqref{eq:armijo} fails with $\alpha=\alpha^{(i)}$, set
$\alpha^{(i+1)}:=\tau\alpha^{(i)}$ and $i:=i+1$. Set $\alpha_{k}:=\alpha^{(i)}$.
\end{algorithm}

\begin{lemma}[Range of admissible stepsizes]\label{lem:armijo-range}
Let $f\in\mathcal{C}^{1}$ with $\grad f$ Lipschitz continuous with constant $L$.
At iteration $k$ of GLM, the Armijo condition \eqref{eq:armijo} is satisfied
for all $\alpha\in[0,\alpha_{k,\max}]$, where
\[
\alpha_{k,\max}=\frac{(\beta-1)\,\grad f(x^{k})\trans s^{k}}{L\norm{s^{k}}^{2}}.
\]
\end{lemma}
\begin{proof}
  First note that $\alpha_{k,\max} > 0$ as $\beta < 1$ and $s^{k}$ is a descent direction. Let $s^{k}$ be a descent direction and let $\alpha \in [0, \alpha_{k,\max}]$. By Taylor's theorem, there exists $\tilde{\alpha} \in (0, \alpha)$ such that
  \begin{align*}
    f(x^{k} + \alpha s^{k}) 
    &= f(x^{k}) + \alpha \grad f(x^{k} + \tilde{\alpha} s^{k})\trans s^{k} \\
    &= f(x^{k}) + \alpha \grad f(x^{k})\trans s^{k} + \alpha \bigl[\grad f(x^{k} + \tilde{\alpha} s^{k}) - \grad f(x^{k})\bigr]\trans s^{k} \\
    &\leq f(x^{k}) + \alpha \grad f(x^{k})\trans s^{k} + \alpha L\norm{\tilde{\alpha}s^{k}} \cdot \norm{s^{k}} \\
    &\leq f(x^{k}) + \alpha \grad f(x^{k})\trans s^{k} + \alpha^2 L\norm{s^{k}}^2 \\
    &\leq f(x^{k}) + \alpha \grad f(x^{k})\trans s^{k} + \alpha^2(\beta - 1)\grad f(x^{k})\trans s^{k}  \\
    &\leq f(x^{k}) + \beta\alpha\grad f(x^{k})\trans s^{k},
  \end{align*}
  which satisfies the Armijo condition.
\end{proof}

\begin{lemma}[Lower bound on bArmijo stepsize]\label{lem:bArmijo-step}
Under the conditions of Lemma~\ref{lem:armijo-range}, the bArmijo stepsize
$\alpha_{k}$ satisfies $\alpha_{k}\ge\min\{\alpha^{(0)},\,\tau\alpha_{k,\max}\}$
for all $k\ge 0$.
\end{lemma}
\begin{proof}
  Note that $\alpha_{k} = \alpha^{(0)}$ if bArmijo terminates at the start. Suppose this is not the case. By Lemma~\ref{lem:armijo-range}, bArmijo terminates whenever $\alpha \leq \alpha_{k,\max}$. Let $i - 1$ be the last iteration before termination, so $\alpha^{(i - 1)} > \alpha_{k,\max}$ and $\tau\alpha^{(i - 1)} = \alpha^{(i)} \leq \alpha_{k,\max}$. But then
  \[
    \alpha_{k} = \alpha^{(i)} = \tau\alpha^{(i - 1)} > \tau\alpha_{k,\max}.
  \]
\end{proof}

\section{Global convergence of GLM}

\begin{theorem}[Global convergence of GLM with bArmijo]\label{thm:GLM-global}
Let $f\in\mathcal{C}^{1}(\R^{n})$ be bounded below by $f_{\mathrm{low}}$, and
$\grad f$ be Lipschitz continuous. Apply GLM with bArmijo linesearch
($\eps:=0$). Then either there exists $\ell\ge 0$ with $\grad f(x^{\ell})=0$,
or
\[
\lim_{k\to\infty}\min\!\left\{\frac{|\grad f(x^{k})\trans s^{k}|}{\norm{s^{k}}},\,
   |\grad f(x^{k})\trans s^{k}|\right\}=0.
\]
\end{theorem}
\begin{proof}
  Suppose $\grad f(x^{k}) \neq 0$ for all $k \geq 0$. Then Armijo condition with $x^{k+1} = x^{k} + \alpha_{k} s^{k}$ gives
  \[
    f(x^{k+1}) \leq f(x^{k}) + \beta \alpha_{k} \grad f(x^{k})\trans s^{k} \implies f(x^{k}) - f(x^{k+1}) \geq \beta \alpha_{k} |\grad f(x^{k})\trans s^{k}|.
  \]
  for all $k \geq 0$. By the telescoping sum,
  \[
    f(x^{0}) - f(x^{k+1}) \geq \sum_{i = 0}^{k} \beta \alpha_{i} |\grad f(x^{i})\trans s^{i}|.
  \]
  Since $f$ is bounded below by $f_{\mathrm{low}}$, 
  \[
    \infty > f(x^{0}) - f_{\mathrm{low}} \geq \sum_{i = 0}^{\infty} \beta \alpha_{i} |\grad f(x^{i})\trans s^{i}|.
  \]
  Thus, $\alpha_{k} |\grad f(x^{k})\trans s^{k}| \to 0$ as $k \to \infty$. Recall from Lemma~\ref{lem:bArmijo-step} that $\alpha_{k} \geq \min\{\alpha^{(0)}, \tau\alpha_{k,\max}\}$ for all $k \geq 0$. Define
  \[
    \mathcal{K}_1 \coloneqq \{k : \alpha^{(0)} \geq \tau\alpha_{k,\max}\}, \qquad \mathcal{K}_2 \coloneqq \{k : \alpha^{(0)} < \tau\alpha_{k,\max}\}.
  \]
  By Lemma~\ref{lem:bArmijo-step}, for $k \in \mathcal{K}_1$,
  \[
    \alpha_{k} |\grad f(x^{k})\trans s^{k}| \geq \frac{(1 - \beta)\tau}{L} \cdot \left(\frac{|\grad f(x^{k})\trans s^{k}|}{\norm{s^{k}}}\right)^2 \geq 0.
  \]
  But then $\alpha_{k} |\grad f(x^{k})\trans s^{k}| \to 0$ as $k \to \infty$ in $\mathcal{K}_1$, so
  \[
    \lim_{\substack{k \to \infty \\ k \in \mathcal{K}_1}} \frac{|\grad f(x^{k})\trans s^{k}|}{\norm{s^{k}}} \to 0.
  \]
  On the other hand, by Lemma~\ref{lem:bArmijo-step}, for $k \in \mathcal{K}_2$ we have $\alpha_{k} \geq \alpha^{(0)}$, so
  \[
    0 \leq \lim_{\substack{k \to \infty \\ k \in \mathcal{K}_2}} |\grad f(x^{k})\trans s^{k}| \leq \frac{1}{\alpha^{(0)}} \cdot \lim_{\substack{k \to \infty \\ k \in \mathcal{K}_2}} \alpha_{k} |\grad f(x^{k})\trans s^{k}| \to 0.
  \]
  Combining both cases, we have
  \[
    \lim_{k\to\infty}\min\!\left\{\frac{|\grad f(x^{k})\trans s^{k}|}{\norm{s^{k}}},\, |\grad f(x^{k})\trans s^{k}|\right\}=0.
  \]
\end{proof}

\begin{remark}\label{rem:angle}
Writing $\cos\theta_{k}=\frac{-\grad f(x^{k})\trans s^{k}}{\norm{\grad f(x^{k})}\norm{s^{k}}}$,
Theorem~\ref{thm:GLM-global} gives
$\norm{\grad f(x^{k})}\cos\theta_{k}\min\{1,\norm{s^{k}}\}\to 0$.
Hence to ensure $\norm{\grad f(x^{k})}\to 0$ it is not enough that $s^{k}$ be
descent: the angle condition $\cos\theta_{k}\ge\delta>0$ for some $\delta$ is
required.
\end{remark}


%==========================================================
\chapter{The Steepest Descent Method}
%==========================================================

The steepest descent (SD) method is the simplest concrete instance of GLM, obtained
by choosing $s^{k}=-\grad f(x^{k})$. It is a \emph{first-order method}: only
gradient information is used.

\section{The steepest descent direction}

\begin{definition}[Steepest descent direction]
The \emph{steepest descent (SD) direction} at $x^{k}$ is
$s^{k}_{\mathrm{SD}}:=-\grad f(x^{k})$.
\end{definition}

\begin{lemma}[Optimality of $-\grad f$]\label{lem:SD}
The steepest descent direction is the unique global solution of
\[
\min_{s\in\R^{n}}\; f(x^{k})+s\trans\grad f(x^{k})
\quad\text{subject to}\quad \norm{s}=\norm{\grad f(x^{k})}.
\]
\end{lemma}
\begin{proof}
  Note that
  \[
    \argmin_{s\in\R^{n}} f(x^{k}) + s\trans\grad f(x^{k}) = \argmin_{s\in\R^{n}} s\trans\grad f(x^{k}).
  \]
  By Cauchy-Schwarz inequality,
  \[
    |s\trans\grad f(x^{k})| \leq \norm{s}\norm{\grad f(x^{k})} = \norm{\grad f(x^{k})}^2,
  \]
  with equality if and only if $s = \pm\grad f(x^{k})$. In particular,
  \[
    s\trans\grad f(x^{k}) \geq -\norm{\grad f(x^{k})}^2,
  \]
  which is only achieved when $s = -\grad f(x^{k})$.
\end{proof}

\section{Global convergence of steepest descent}

\begin{theorem}[Global convergence of SD with bArmijo]\label{thm:SD-global}
Let $f\in\mathcal{C}^{1}(\R^{n})$ be bounded below with $\grad f$ Lipschitz
continuous. Apply GLM with $s^{k}=-\grad f(x^{k})$ and bArmijo linesearch.
Then either there exists $\ell$ with $\grad f(x^{\ell})=0$, or
$\norm{\grad f(x^{k})}\to 0$ as $k\to\infty$.
\end{theorem}
\begin{proof}
  Suppose there are no $\ell$ with $\grad f(x^{\ell}) = 0$. Then by Theorem~\ref{thm:GLM-global}, 
  \[
  \lim_{k\to\infty}\min \left\{\frac{\norm{\grad f(x^{k})}^2}{\norm{-\grad f(x^{k})}}, \norm{\grad f(x^{k})}\right\} = \norm{\grad f(x^{k})} = 0.
  \]
\end{proof}

\section{Local rate of convergence}

\begin{definition}[Rates of convergence]
A sequence $\{x^{k}\}\subset\R^{n}$ converges to $x^{*}$ with rate $p\ge 1$ if
there exist $\rho>0$ and $k_{0}$ with
$\norm{x^{k+1}-x^{*}}\le\rho\norm{x^{k}-x^{*}}^{p}$ for all $k\ge k_{0}$.
The cases $p=1$ ($\rho<1$) and $p=2$ are called \emph{linear} and
\emph{quadratic} convergence respectively. The convergence is
\emph{superlinear} if $\norm{x^{k+1}-x^{*}}/\norm{x^{k}-x^{*}}\to 0$.
\end{definition}

\begin{theorem}[Local linear convergence of SD]\label{thm:SD-local}
Let $f\in\mathcal{C}^{2}$ and let $x^{*}$ be a local minimizer of $f$ with
$\hess f(x^{*})\succ 0$ and eigenvalues $0<\lambda^{*}_{\min}\le\lambda^{*}_{\max}$.
Apply SD with exact linesearch. If $x^{k}\to x^{*}$, then the convergence is
linear with factor
\[
\rho\le \frac{\kappa(x^{*})-1}{\kappa(x^{*})+1},
\qquad
\kappa(x^{*})=\frac{\lambda^{*}_{\max}}{\lambda^{*}_{\min}}.
\]
\qed
\end{theorem}

\section{Summary of steepest descent}

\begin{itemize}[leftmargin=2em,itemsep=0pt]
  \item \textbf{Inexpensive:} only first derivatives required.
  \item \textbf{Strong global convergence} under weak assumptions, with second-order
        optimality guarantees for the generated solution (see e.g.\ saddle-point
        avoidance properties in Lecture 5).
  \item \textbf{Scale-dependent:} a poor scaling of variables produces a
        steepest-descent direction that gives little progress. Newton's method
        does not have this defect (Theorem~\ref{thm:newton-scale}).
  \item \textbf{Slow asymptotic convergence}: when $\kappa(x^{*})$ is large,
        $\rho_{SD}=(\kappa-1)/(\kappa+1)$ is very close to $1$ and the iterates
        ``zig-zag''. SD is therefore not used in general-purpose software, but it
        remains useful for certain very well-conditioned problems with structure
        (e.g.\ compressed sensing).
\end{itemize}


%==========================================================
\chapter{Newton's Method}
\label{ch:newton}
%==========================================================

Newton's method for unconstrained minimization can be derived in two equivalent
ways:
\begin{enumerate}[leftmargin=2em,itemsep=0pt]
  \item \emph{Quadratic model.} The Newton step $s^{k}$ minimizes the second-order
        Taylor approximation of $f$ at $x^{k}$ (when the Hessian is positive
        definite). Compare this with steepest descent, whose direction minimizes
        the first-order Taylor approximation under a norm constraint.
  \item \emph{Newton's root-finding.} Stationarity is the equation $\grad f(x)=0$,
        an $n\times n$ nonlinear system. Applying Newton's root-finding scheme to
        $r(x):=\grad f(x)$ produces
        $x^{k+1}=x^{k}-[J(x^{k})]^{-1}r(x^{k})=x^{k}-[\hess f(x^{k})]^{-1}\grad f(x^{k})$,
        recovering the Newton iterate.
\end{enumerate}

\section{The Newton direction}

\begin{definition}[Newton direction]
For $f\in\mathcal{C}^{2}$ with $\hess f(x^{k})$ nonsingular, the
\emph{Newton direction} at $x^{k}$ is the solution $s^{k}$ of
\begin{equation}\label{eq:newton-system}
\hess f(x^{k})\,s^{k}=-\grad f(x^{k}).
\end{equation}
The \emph{(pure) Newton iteration} sets $x^{k+1}:=x^{k}+s^{k}$.
\end{definition}

\begin{lemma}[Descent property of the Newton direction]\label{lem:newton-descent}
If $\hess f(x^{k})\succ 0$ and $\grad f(x^{k})\ne0$, then the Newton direction is
a descent direction. Moreover, $s^{k}$ is the unique minimizer of the
second-order Taylor model
$q_{k}(s):=f(x^{k})+s\trans\grad f(x^{k})+\tfrac12 s\trans\hess f(x^{k})s$.
\end{lemma}
\begin{proof}
  Put $H_{k}:=\hess f(x^{k})$ and $g_{k}:=\grad f(x^{k})$. Since
  $H_{k}\succ 0$, it is nonsingular and the Newton equation gives
  \[
    s^{k} = -H_{k}^{-1}g_{k}.
  \]
  Also $H_{k}^{-1}\succ 0$, so, as $g_{k}\ne0$,
  \[
    g_{k}\trans s^{k} = -g_{k}\trans H_{k}^{-1}g_{k}<0,
  \]
  which makes $s^{k}$ a descent direction.

  It remains to prove the model-minimizing property. Completing the square, we may write
  \[
    q_{k}(s) = f(x^{k})-\frac12 g_{k}\trans H_{k}^{-1}g_{k} + \frac12 (s+H_{k}^{-1}g_{k})\trans H_{k}(s+H_{k}^{-1}g_{k}).
  \]
  The last term is nonnegative, and it vanishes if and only if
  $s=-H_{k}^{-1}g_{k}=s^{k}$. Hence $s^{k}$ is the unique global minimizer of
  $q_{k}$.
\end{proof}

\section{Local convergence of pure Newton}

\begin{theorem}[Local quadratic convergence of pure Newton]\label{thm:newton-local}
Let $f\in\mathcal{C}^{2}(\R^{n})$, $\grad f(x^{*})=0$, $\hess f(x^{*})$ nonsingular,
and $\hess f$ locally Lipschitz continuous at $x^{*}$. If $x^{k_{0}}$ is sufficiently
close to $x^{*}$ for some $k_{0}\ge 0$, then $x^{k}$ is well-defined for all
$k\ge k_{0}$ and $x^{k}\to x^{*}$ at a quadratic rate.
\end{theorem}
\begin{proof}
  Write $g(x):=\grad f(x)$, $H(x):=\hess f(x)$, and
  $e^{k}:=x^{k}-x^{*}$. Since $H(x^{*})$ is nonsingular and $H$ is continuous,
  there is a radius $\rho>0$ such that $H(x)$ is nonsingular for
  $x\in N(x^{*},\rho)$. Shrinking $\rho$ if necessary, local Lipschitz
  continuity gives a constant $L>0$ such that
  \[
    \norm{H(y)-H(x)}\le L\norm{y-x}
    \quad\text{for all }x,y\in N(x^{*},\rho),
  \]
  and continuity of $x\mapsto H(x)^{-1}$ gives a constant $M>0$ such that
  \[
    \norm{H(x)^{-1}}\le M
    \quad\text{for all }x\in N(x^{*},\rho).
  \]
  For any $x\in N(x^{*},\rho)$, Taylor's theorem for $g$ with Lipschitz derivative gives
  \[
    g(x^{*})=g(x)+H(x)(x^{*}-x)+r(x),
    \qquad
    \norm{r(x)}\le \frac{L}{2}\norm{x-x^{*}}^{2}.
  \]
  Since $g(x^{*})=0$, the pure Newton iterate
  $N(x):=x-H(x)^{-1}g(x)$ satisfies
  \[
    N(x)-x^{*}
    =H(x)^{-1}\bigl(H(x)(x-x^{*})-g(x)\bigr)
    =H(x)^{-1}r(x).
  \]
  Hence
  \[
    \norm{N(x)-x^{*}}
    \le \frac{ML}{2}\norm{x-x^{*}}^{2}.
  \]

  Choose $\delta\in(0,\rho)$ so small that
  $\frac{ML}{2}\delta\le\frac12$. If $x^{k_{0}}\in N(x^{*},\delta)$, the above
  estimate shows inductively that every Newton iterate is well-defined and
  remains in $N(x^{*},\delta)$, since
  \[
    \norm{x^{k+1}-x^{*}}
    \le \frac{ML}{2}\norm{x^{k}-x^{*}}^{2}
    \le \frac12\norm{x^{k}-x^{*}}.
  \]
  Therefore $x^{k}\to x^{*}$, and the sharper estimate
  \[
    \norm{x^{k+1}-x^{*}}
    \le \frac{ML}{2}\norm{x^{k}-x^{*}}^{2}
  \]
  gives quadratic convergence.
\end{proof}

\begin{theorem}[Scale invariance of Newton]\label{thm:newton-scale}
Newton's method (with or without linesearch) is scale invariant under linear
transformations of the variables. Explicitly, if $A\in\R^{n\times n}$ is
nonsingular, $y=Ax$, and $\bar f(y):=f(A^{-1}y)$, then the Newton direction $s_{y}$
for $\bar f$ at $y$ satisfies $y+\alpha s_{y}=A(x+\alpha s_{x})$.
\end{theorem}

\begin{proof}
  Note that
  \[
    \grad \bar f(y) = (A^{-1})\trans \grad f(x), \quad \hess \bar f(y) = (A^{-1})\trans \hess f(x) A^{-1}.
  \]
  Thus,
  \begin{align*}
    s_y = -\bigl[\hess \bar f(y)\bigr]^{-1} \grad \bar f(y) = \bigl(A \bigl[\hess f(x)\bigr]^{-1} A\trans\bigr) \bigl((A^{-1})\trans \grad f(x)\bigr) = A s_x.
  \end{align*}
\end{proof}

\begin{remark}[Disadvantages of Newton's method]
Despite its fast local convergence, the pure Newton iteration suffers several
serious drawbacks:
\begin{itemize}[leftmargin=2em,itemsep=0pt]
  \item $s^{k}$ is not well-defined when $\hess f(x^{k})$ is singular.
  \item $s^{k}$ may not be a descent direction when $\hess f(x^{k})$ is not
        positive definite.
  \item Pure Newton may be attracted to local maxima or saddle points (the local
        convergence Theorem~\ref{thm:newton-local} only requires $\hess f(x^{*})$
        nonsingular, not positive definite).
  \item Pure Newton may fail to converge if $x^{0}$ is far from a solution; for
        example, applied to $f(x)=-x^{6}/6+x^{4}/4+2x^{2}$ from $x^{0}=1$, it
        cycles between $\pm 1$ (neither stationary).
\end{itemize}
This is why a linesearch (or trust-region) safeguard is needed to make Newton's
method globally convergent.
\end{remark}

\section{Linesearch Newton (damped Newton) method}

\begin{algorithm}[Linesearch Newton method]
Choose $\eps>0$ and $x^{0}\in\R^{n}$. While $\norm{\grad f(x^{k})}>\eps$:
solve $\hess f(x^{k})s^{k}=-\grad f(x^{k})$, then set $x^{k+1}=x^{k}+\alpha_{k}s^{k}$
with $\alpha_{k}\in(0,1]$ chosen by, e.g., bArmijo linesearch.
\end{algorithm}

\begin{theorem}[Local convergence of linesearch Newton]\label{thm:newton-armijo}
Let $f\in\mathcal{C}^{2}(\R^{n})$ with $\hess f$ Lipschitz continuous and positive
definite at the iterates. Apply Newton's method with bArmijo linesearch using
$\beta<1/2$ and $\alpha^{(0)}=1$. If $x^{k}\to x^{*}$ with $\grad f(x^{*})=0$ and
$\hess f(x^{*})\succ 0$, then $\alpha_{k}=1$ for all $k$ sufficiently large and
the convergence is asymptotically quadratic.
\end{theorem}
\begin{proof}
  Put $g_{k}:=\grad f(x^{k})$, $H_{k}:=\hess f(x^{k})$, and let
  $s^{k}=-H_{k}^{-1}g_{k}$ be the Newton direction. Since
  $x^{k}\to x^{*}$ and $H(x^{*})\succ0$, there exist $m>0$ and $k_{0}$ such that
  \[
    H_{k}\succeq m I
    \quad\text{for all }k\ge k_{0}.
  \]
  In particular, for $k\ge k_{0}$,
  \[
    g_{k}\trans s^{k}=-(s^{k})\trans H_{k}s^{k}
    \le -m\norm{s^{k}}^{2}.
  \]
  Also $g_{k}\to0$ and $H_{k}^{-1}$ is bounded for large $k$, so
  $s^{k}\to0$.

  Let $L$ be a Lipschitz constant for $\hess f$ on a neighbourhood of $x^{*}$
  containing $x^{k}$ and $x^{k}+s^{k}$ for all sufficiently large $k$. Taylor's
  theorem with Lipschitz Hessian gives
  \[
    f(x^{k}+s^{k})
    \le f(x^{k})+g_{k}\trans s^{k}
      +\frac12 (s^{k})\trans H_{k}s^{k}
      +\frac{L}{6}\norm{s^{k}}^{3}.
  \]
  Using $H_{k}s^{k}=-g_{k}$, this becomes
  \[
    f(x^{k}+s^{k})
    \le f(x^{k})+\frac12 g_{k}\trans s^{k}
      +\frac{L}{6}\norm{s^{k}}^{3}.
  \]
  Since $\beta<1/2$ and $s^{k}\to0$, for all sufficiently large $k$,
  \[
    \frac{L}{6}\norm{s^{k}}
    \le \left(\frac12-\beta\right)m.
  \]
  Therefore
  \[
    \frac{L}{6}\norm{s^{k}}^{3}
    \le \left(\frac12-\beta\right)(s^{k})\trans H_{k}s^{k}
    = \left(\frac12-\beta\right)(-g_{k}\trans s^{k}),
  \]
  and hence
  \[
    f(x^{k}+s^{k})
    \le f(x^{k})+\beta\, g_{k}\trans s^{k}.
  \]
  Thus the full step $\alpha=1$ satisfies the Armijo condition. Since bArmijo
  starts with $\alpha^{(0)}=1$, it accepts $\alpha_{k}=1$ for all sufficiently
  large $k$.

  From that point on the method is pure Newton. The hypotheses of
  Theorem~\ref{thm:newton-local} hold at $x^{*}$, so the tail of the sequence
  satisfies
  \[
    \norm{x^{k+1}-x^{*}}\le C\norm{x^{k}-x^{*}}^{2}
  \]
  for some constant $C>0$ and all sufficiently large $k$. This is asymptotic
  quadratic convergence.
\end{proof}

\section{Global convergence with general second-order models}

In GLM, let $B_{k}$ be symmetric positive definite for all $k$, and let $s^{k}$
solve
\begin{equation}\label{eq:Bksk}
B_{k}s^{k}=-\grad f(x^{k}).
\end{equation}

\begin{theorem}[Global convergence with bounded $B_{k}$]\label{thm:Bk-global}
Let $f\in\mathcal{C}^{1}(\R^{n})$ be bounded below with $\grad f$ Lipschitz
continuous. Apply GLM with $s^{k}$ from \eqref{eq:Bksk} and bArmijo linesearch
($\eps:=0$). Suppose the eigenvalues of $B_{k}$ are uniformly bounded above and
uniformly bounded away from zero, independently of $k$. Then either there exists
$\ell$ with $\grad f(x^{\ell})=0$, or $\norm{\grad f(x^{k})}\to 0$ as $k\to\infty$.
\end{theorem}
\begin{proof}
  Write $H_k = \hess f(x^{k})$ and $g_k = \grad f(x^{k})$. Note that $B_k$ is nonsingular, so we may write $s^{k} = -B_k^{-1}g_k$. Assume that $g_k \neq 0$ for all $k \geq 0$. Then by Theorem~\ref{thm:GLM-global}
  \begin{equation}\label{eq:Bk-global-condition}
    \lim_{k\to\infty}\min\!\left\{\frac{|g_k\trans s^{k}|}{\norm{s^{k}}}, |g_k\trans s^{k}|\right\}=0.
  \end{equation}
  We are left to show that this implies that $\norm{\grad f(x^{k})} \to 0$ as $k \to \infty$.

  Since $B_k$ is symmetric, $B_k$ admits real eigenvalues by the spectral theorem. Since the eigenvalues are uniformly bounded above and away from zero, there exists constants $\lambda_{\max}$ and $\lambda_{\min}$ such that $\lambda_{\max} \geq \lambda_{\max}(B_k) \geq \lambda_{\min}(B_k) \geq \lambda_{\min}$ for all $k \geq 0$. In particular, for any $s \in \R^n$ and any $B_k$,
  \[
    \lambda_{\min} \norm{s}^2 \leq \lambda_{\min}(B_k) \norm{s}^2 \leq s\trans B_k s \leq \lambda_{\max}(B_k) \norm{s}^2 \leq \lambda_{\max} \norm{s}^2 \implies \lambda_{\min} \leq \frac{s\trans B_k s}{\norm{s}^2} \leq \lambda_{\max}.
  \]
  Since $\grad f$ is Lipschitz continuous, $\norm{H_k}$ is uniformly bounded above. Thus, by possibly enlarging $\lambda_{\max}$, we have $\lambda_{\max}(H_k) \leq \lambda_{\max}$ for all $k \geq 0$. 

  By definition of $s^{k}$ and $g_k \neq 0$, 
  \[
    |g_k\trans s^{k}| = |g_k\trans B_k^{-1}g_k| \geq \lambda_{\min}(B_k^{-1}) \norm{g_k}^2 \geq \frac{\norm{g_k}^2}{\lambda_{\max}}.
  \]
  On the other hand,
  \[
    \norm{s^{k}}^2 = g_k\trans B_k^{-2}g_k \leq \lambda_{\max}(B_k^{-2})\norm{g_k}^2 \leq \frac{\norm{g_k}^2}{\lambda_{\min}^2} \implies \frac{1}{\norm{s^{k}}} \geq \frac{\lambda_{\min}}{\norm{g_k}}.
  \]
  Plugging into \eqref{eq:Bk-global-condition}, we have
  \[
    \lim_{k\to\infty} \min\!\left\{\frac{\lambda_{\min}}{\lambda_{\max}} \norm{g_k}, \frac{\norm{g_k}^2}{\lambda_{\max}}\right\} \leq \lim_{k\to\infty}\min\!\left\{\frac{|g_k\trans s^{k}|}{\norm{s^{k}}}, |g_k\trans s^{k}|\right\} \to 0 \implies \lim_{k \to \infty} g_k = 0.
  \]
\end{proof}


\chapter{Quasi-Newton Methods and Nonlinear Least-Squares}
\label{ch:qn-nls}

Newton's method has fast local convergence but each iteration is expensive: forming
$\hess f(x^{k})$ costs $O(n^{2})$ second-derivative evaluations, and solving the
linear system costs $O(n^{3})$. \emph{Quasi-Newton methods} aim for a sweet spot:
faster than steepest descent, but cheaper per iteration than Newton, by using
\emph{secant approximations} $B_{k}\approx\hess f(x^{k})$ built from gradient
information already computed.

The design principles for $B_{k+1}$ form a ``wish list'':
\begin{itemize}[leftmargin=2em,itemsep=0pt]
  \item built from already-computed quantities $\grad f(x^{k+1}),\grad f(x^{k}),
        \dots,B_{k},s^{k}$;
  \item symmetric, nonsingular, and ideally positive definite;
  \item ``close'' to $B_{k}$ and a ``cheap'' update of $B_{k}$;
  \item asymptotically, $B_{k}\to\hess f(x^{*})$.
\end{itemize}
No single update satisfies the entire list; different methods make different
trade-offs.

\begin{remark}[Approximating the Hessian by finite differences]
An alternative when second derivatives are unavailable is to approximate
$[\hess f(x)]e_{i}\approx \frac{1}{h}[\grad f(x+he_{i})-\grad f(x)]$ for
$i=1,\dots,n$, costing $n+1$ gradient evaluations. Choice of $h$ is delicate (too
large $\to$ inaccurate; too small $\to$ cancellation errors).
\end{remark}

\section{The secant equation}

In quasi-Newton methods, $\hess f(x^{k+1})$ is approximated by a symmetric matrix
$B_{k+1}$ satisfying the \emph{secant equation}
\begin{equation}\label{eq:secant}
B_{k+1}\,\delta^{k}=\gamma^{k},
\qquad \delta^{k}:=x^{k+1}-x^{k},
\qquad \gamma^{k}:=\grad f(x^{k+1})-\grad f(x^{k}).
\end{equation}

\begin{lemma}[Quadratic case]\label{lem:secant-quad}
If $f(x)=g\trans x+\tfrac12 x\trans Hx$, then \eqref{eq:secant} is satisfied with
$B_{k+1}=\hess f=H$.
\end{lemma}
\begin{proof}
  Since $\grad f(x) = Hx + g$,
  \[
    \grad f(x^{k+1}) - \grad f(x^{k}) = H(x^{k+1} - x^{k}),
  \]
  which satisfies the secant equation.
\end{proof}

\section{The BFGS update}

There are infinitely many ways to satisfy the secant equation. One choice is the
\emph{symmetric rank-1 (SR1)} update $B_{k+1}=B_{k}+u^{k}(u^{k})\trans$, which is
cheap ($O(n^{2})$ per iteration via the Sherman--Morrison--Woodbury formula) but
may fail to be positive definite, leading to non-descent directions.
The \emph{BFGS} update (Broyden, Fletcher, Goldfarb, Shanno, independently 1970)
is a rank-2 update that preserves positive definiteness.

\begin{definition}[BFGS update]
The \emph{BFGS} update is the rank-2 symmetric update
\begin{equation}\label{eq:BFGS}
B_{k+1}=B_{k}-\frac{B_{k}\delta^{k}(\delta^{k})\trans B_{k}}{(\delta^{k})\trans B_{k}\delta^{k}}
+\frac{\gamma^{k}(\gamma^{k})\trans}{(\delta^{k})\trans\gamma^{k}}.
\end{equation}
\end{definition}

\begin{theorem}[Properties of the BFGS update]\label{thm:BFGS-pd}
If $B_{k}$ is symmetric positive definite and $(\delta^{k})\trans\gamma^{k}>0$,
then the BFGS update \eqref{eq:BFGS} is well-defined, symmetric and positive
definite, and satisfies the secant equation \eqref{eq:secant}.
\end{theorem}
\begin{proof}
  It is easy to see that $B_{k + 1}$ is symmetric, as the two terms after $B_k$ are symmetric. For $x \in \R^n$ with $x \neq 0$, 
  \[
    x\trans B_{k+1} x = x\trans B_{k}x - \frac{\bigl[(\delta^{k})\trans B_{k} x\bigr]^2}{(\delta^{k})\trans B_{k}\delta^{k}}
    +\frac{\bigl[(\gamma^{k})\trans x\bigr]^2}{(\delta^{k})\trans\gamma^{k}}.
  \]
  Since $B_k$ is positive definite, we may define the inner product $\langle x, y \rangle_B = x\trans B_k y$ induced by $B_k$. Since $(\delta^{k})\trans\gamma^{k}>0$, the term $\frac{[(\gamma^{k})\trans x]^2}{(\delta^{k})\trans\gamma^{k}}$ is positive, and so by Cauchy-Schwarz inequality,
  \[
    x\trans B_{k+1} x > \langle x, x \rangle_B - \frac{\langle \delta^k, x \rangle_B^2}{\langle \delta^k, \delta^k \rangle_B} \geq \langle x, x \rangle_B - \frac{\langle x, x \rangle_B \cdot \langle \delta^k, \delta^k \rangle_B}{\langle \delta^k, \delta^k \rangle_B} = 0.
  \]
  Since
  \[
    B_{k + 1}\delta^k = B_k\delta^k - \frac{B_k\delta^k[(\delta^k)\trans B_k\delta^k]}{(\delta^k)\trans B_k\delta^k} + \frac{\gamma^k[(\gamma^k)\trans\delta^k]}{(\delta^k)\trans\gamma^k} = \gamma^k,
  \]
  the secant equation is satisfied.
\end{proof}

\begin{remark}
The condition $(\delta^{k})\trans\gamma^{k}>0$ is guaranteed by use of the
Wolfe linesearch.
\end{remark}

\begin{theorem}[Convergence of BFGS]\label{thm:BFGS-conv}
Under suitable conditions on $f$ and using the Wolfe linesearch, the BFGS method
is globally convergent and locally Q-superlinearly convergent. \qed
\end{theorem}

\begin{remark}[BFGS in practice]
BFGS satisfies essentially all the items on the wish list: it is symmetric and
positive definite (so $s^{k}=-B_{k}^{-1}\grad f(x^{k})$ is descent), close to $B_{k}$,
cheap to update (in practice the Cholesky factors of $B_{k}$ are stored and updated
directly). Applied with exact linesearch to a strictly convex quadratic, $B_{k}$
converges to $\hess f$ in at most $n$ iterations. BFGS suffers much less from poor
scaling than steepest descent. It has been the most popular general-purpose method
when second derivatives are unavailable, and is the basis of MATLAB's
\texttt{fminunc}.
\end{remark}

\section{Nonlinear least-squares}

Many problems (model fitting, parameter estimation, data assimilation) take the
form: choose $x$ to minimise the sum of squared residuals $r_{j}(x)^{2}$. The
problem is called \emph{zero-residual} if $r(x^{*})=0$ (the model fits exactly)
and \emph{nonzero-residual} otherwise. The special structure of $f$ allows
significant savings over generic Newton methods.

The nonlinear least-squares problem is
\begin{equation}\label{eq:NLS}
\min_{x\in\R^{n}}\; f(x)=\tfrac12\norm{r(x)}^{2}=\tfrac12\sum_{j=1}^{m}r_{j}(x)^{2},
\end{equation}
where $r:\R^{n}\to\R^{m}$ is smooth with Jacobian $J(x)$.

\begin{lemma}[Derivatives of $f$]\label{lem:NLS-deriv}
For $f$ in \eqref{eq:NLS},
\[
\grad f(x)=J(x)\trans r(x),\qquad
\hess f(x)=J(x)\trans J(x)+\sum_{j=1}^{m}r_{j}(x)\hess r_{j}(x).
\]
\end{lemma}
\begin{proof}
  \[
    [\grad f(x)]_i = \frac{\partial f}{\partial x_i}(x) = \sum_{j = 1}^m r_j(x) \cdot \frac{\partial r_j}{\partial x_i}(x) = r(x)\trans \frac{\partial r}{\partial x_i}(x) = r(x)\trans J(x)e_i.
  \]
  \begin{align*}
    [\hess f(x)]_{ik} 
    &= \frac{\partial f}{\partial x_i x_k}(x) \\
    &= \sum_{j = 1}^m \frac{\partial}{\partial x_k} \left[r_j(x) \cdot \frac{\partial r_j}{\partial x_i}(x)\right] \\
    &= \sum_{j = 1}^m \frac{\partial r_j}{\partial x_i}(x) \cdot \frac{\partial r_j}{\partial x_k}(x) + r_j(x) \cdot \frac{\partial^2 r_j}{\partial x_i \partial x_k}(x) \\
    &= [J^T(x)J(x)]_{ik} + \left[\sum_{j=1}^{m}r_{j}(x)\hess r_{j}(x)\right]_{ik}.
  \end{align*}
\end{proof}

\section{The Gauss-Newton method}

\paragraph{Motivation.} The Hessian decomposes as $\hess f(x)=J(x)\trans J(x)
+\sum_{j}r_{j}(x)\hess r_{j}(x)$. When residuals are small ($r_{j}(x^{*})\approx 0$)
or the model is nearly linear ($\hess r_{j}\approx 0$), the second term is
negligible and $\hess f(x)\approx J(x)\trans J(x)$. The Gauss-Newton method
substitutes this approximation into the Newton system, dispensing with second
derivatives entirely.

\begin{definition}[Gauss-Newton direction]
The \emph{Gauss-Newton (GN) direction} solves the linear system
\begin{equation}\label{eq:GN}
J(x^{k})\trans J(x^{k})\,s^{k}=-J(x^{k})\trans r(x^{k}).
\end{equation}
\end{definition}

\begin{lemma}[Descent property]\label{lem:GN-descent}
If $J(x^{k})$ has full column rank, then the Gauss-Newton direction is a descent
direction for $f$.
\end{lemma}
\begin{proof}
  Since $J(x^k)$ has full column rank, $J(x^k)\trans J(x^k)$ is positive definite, and so 
  \[
    [\grad f(x^k)]\trans s^k = -r(x^k)\trans J(x^k)  J(x^k)\trans r(x^k) < 0.
  \]
\end{proof}

\begin{theorem}[Convergence of Gauss-Newton]\label{thm:GN-conv}
Apply Gauss-Newton with bArmijo linesearch.
\begin{enumerate}[leftmargin=2em,itemsep=0pt]
\item (Global.) If $J(x^{k})$ is uniformly full rank, $f\in\mathcal{C}^{1}$ bounded below
and $\grad f$ Lipschitz, then $\norm{\grad f(x^{k})}\to 0$.
\item (Local.) If $r(x^{*})=0$ and $J(x^{*})$ has full rank, the convergence is quadratic; if $r(x^{*})\ne 0$ but small, the convergence is linear. \qed
\end{enumerate}
\end{theorem}

\begin{remark}[Newton vs.\ Gauss-Newton]\
\begin{itemize}[leftmargin=2em,itemsep=0pt]
  \item \textbf{Cost per iteration:} Newton $>$ Gauss-Newton (no second
        derivatives needed for GN).
  \item \textbf{Newton direction may be ascent} if $\hess f(x^{k})$ is not positive
        definite. In contrast, the GN direction is always descent provided
        $J(x^{k})$ has full column rank.
  \item \textbf{Local rate:} Newton always quadratic; GN only quadratic for
        zero-residual problems and only linear when $r(x^{*})\ne 0$.
  \item Both methods are unreliable without a linesearch (or trust-region)
        safeguard.
\end{itemize}
\end{remark}


\chapter{Trust-Region Methods}

Trust-region methods are an alternative to linesearch for globalising Newton-type
iterations. The two philosophies differ:
\begin{itemize}[leftmargin=2em,itemsep=0pt]
  \item \textbf{Linesearch:} ``liberal'' in choice of search direction; keeps bad
        behaviour in control by choosing $\alpha_{k}$.
  \item \textbf{Trust-region:} ``conservative'' in choice of step; the full step
        is accepted if a sufficient \emph{model decrease} is realised in the
        objective. Otherwise the step is rejected and the model is refined.
\end{itemize}
A key advantage of TR methods is that they accommodate \emph{indefinite} or
ill-behaved Hessians: the trust-region constraint regularises an otherwise
unbounded model. They also fit naturally with quasi-Newton ideas, since there is
no need for $B_{k}$ to be positive definite.

\section{The trust-region framework}

We approximate $f(x^{k}+s)$ by a local model. Two natural choices are:
\begin{itemize}[leftmargin=2em,itemsep=0pt]
  \item linear: $\ell_{k}(s):=f(x^{k})+s\trans\grad f(x^{k})$;
  \item quadratic: $q_{k}(s):=f(x^{k})+s\trans\grad f(x^{k})+\tfrac12 s\trans
        B_{k} s$, with $B_{k}=\hess f(x^{k})$ in the standard case.
\end{itemize}
\paragraph{Impediments.} The model may not resemble $f$ when $\norm{s}$ is large,
and may be unbounded below: $\ell_{k}$ is always unbounded below (unless
$\grad f(x^{k})=0$), and $q_{k}$ is unbounded below whenever $B_{k}$ is negative
definite or indefinite, and sometimes when $B_{k}$ is only positive
semidefinite. We trust the model only on a region $\norm{s}\le\Delta_{k}$:
At iteration $k$, we form a quadratic model
\begin{equation}\label{eq:TRmodel}
m_{k}(s)=f(x^{k})+s\trans\grad f(x^{k})+\tfrac12 s\trans B_{k}s,
\end{equation}
where $B_{k}$ is symmetric (e.g.\ $B_{k}=\hess f(x^{k})$). The
\emph{trust-region (TR) subproblem} is
\begin{equation}\label{eq:TRsub}\tag{TR}
\min_{s\in\R^{n}}\; m_{k}(s) \quad\text{subject to}\quad \norm{s}\le\Delta_{k},
\end{equation}
with trust-region radius $\Delta_{k}>0$.

\begin{algorithm}[Generic Trust-Region (GTR) method]\label{alg:GTR}
Given $\Delta_{0}>0$, $x^{0}\in\R^{n}$, $\eps>0$. While $\norm{\grad f(x^{k})}>\eps$:
\begin{enumerate}[leftmargin=2em,itemsep=0pt]
  \item Form $m_{k}$ and (approximately) solve \eqref{eq:TRsub} for $s^{k}$ with
        $m_{k}(s^{k})<f(x^{k})$.
  \item Compute $\rho_{k}=\dfrac{f(x^{k})-f(x^{k}+s^{k})}{f(x^{k})-m_{k}(s^{k})}$.
  \item If $\rho_{k}\ge 0.9$: set $x^{k+1}:=x^{k}+s^{k}$, $\Delta_{k+1}:=2\Delta_{k}$.\\
        Else if $\rho_{k}\ge 0.1$: set $x^{k+1}:=x^{k}+s^{k}$, $\Delta_{k+1}:=\Delta_{k}$.\\
        Else: set $x^{k+1}:=x^{k}$, $\Delta_{k+1}:=\tfrac12\Delta_{k}$.
\end{enumerate}
\end{algorithm}

\begin{remark}[The ratio $\rho_{k}$]
$\rho_{k}$ measures the ratio of the \emph{actual function decrease} to the
\emph{predicted model decrease}. $\rho_{k}\approx 1$ indicates the model is
accurate (we expand the trust region to take longer steps); $\rho_{k}\ll 1$
indicates the model is poor (we reject the step and shrink the trust region).
The thresholds $0.1,0.9$ are conventional; other reasonable values appear in
practice.
\end{remark}

\section{The Cauchy point}

The Cauchy point is the trust-region analogue of the steepest descent direction:
it minimises the model $m_{k}$ along the steepest descent direction subject to the
trust-region constraint. Recall (Theorem~\ref{thm:SD-global}) that steepest descent
has strong global convergence properties; insisting that any approximate solution
$s^{k}$ of the TR subproblem do at least as well as the Cauchy point is a
``minimal sufficient decrease'' condition that suffices for global convergence
of GTR.

\begin{definition}[Cauchy point]
The \emph{Cauchy point} of \eqref{eq:TRsub} is $y^{k}_{c}=x^{k}+s^{k}_{c}$, where
$s^{k}_{c}=-\alpha_{k,c}\grad f(x^{k})$ and
\[
\alpha_{k,c}=\argmin_{\alpha>0}\, m_{k}\!\left(-\alpha\grad f(x^{k})\right)
\quad\text{subject to}\quad
\norm{\alpha\grad f(x^{k})}\le\Delta_{k}.
\]
\end{definition}

\begin{lemma}[Cauchy-point decrease]\label{lem:cauchy}
\[
f(x^{k})-m_{k}(s^{k}_{c})\ge \tfrac12\norm{\grad f(x^{k})}\,
\min\!\left\{\Delta_{k},\,\frac{\norm{\grad f(x^{k})}}{\norm{B_{k}}}\right\}.
\]
\end{lemma}
\begin{proof}
  Write $g_k := \grad f(x^{k})$ and $h_{k}=\grad f(x^{k})\trans B_{k}\grad f(x^{k})$. Define
  \[
  \psi(\alpha) := m_k(-\alpha g_k)
              = f(x^{k}) - \alpha\norm{g_k}^{2} + \tfrac{1}{2}\alpha^{2} h_k,
  \qquad \alpha > 0,
  \]
  so that the Cauchy step is $s^{k}_c = -\alpha^{k}_c g_k$ with
  \[
  \alpha^{k}_c \;=\; \arg\min_{0<\alpha\le\bar\alpha}\psi(\alpha),
  \qquad
  \bar\alpha := \frac{\Delta_k}{\norm{g_k}}.
  \]
  Note $\psi'(0) = -\norm{g_k}^{2} < 0$, so $\psi$ is initially decreasing
  and $\alpha^{k}_c > 0$.
  
  \noindent\textbf{Case 1: $h_k \le 0$.}\quad
  Then $\psi$ is concave (or linear) and strictly decreasing on
  $(0,\bar\alpha]$, so the minimum on the trust region is attained at the
  boundary: $\alpha^{k}_c = \bar\alpha = \Delta_k/\norm{g_k}$. Since the
  quadratic term $\tfrac12\alpha^{2}h_k \leq 0$,
  \[
  f(x^{k}) - m_k(s^{k}_c) = \alpha^{k}_c\norm{g_k}^{2} - \tfrac12(\alpha^{k}_c)^{2} h_k \geq \alpha^{k}_c\norm{g_k}^{2} = \Delta_k \norm{g_k}.
  \]
  
  \textbf{Case 2: $h_k > 0$.}\quad
  Now $\psi$ is a strictly convex quadratic in $\alpha$ with unconstrained
  minimiser
  \[
  \alpha_{\min} \;=\; \frac{\norm{g_k}^{2}}{h_k},
  \]
  and $\alpha^{k}_c = \min\{\alpha_{\min},\,\bar\alpha\}$. We consider the
  two sub-cases.
  
  \emph{Sub-case 2a: $\alpha_{\min} \le \bar\alpha$.}\quad
  Then $\alpha^{k}_c = \alpha_{\min}$ and substituting gives
  \[
  f(x^{k}) - m_k(s^{k}_c)
      = \alpha_{\min}\norm{g_k}^{2} - \tfrac12\alpha_{\min}^{2}h_k
      = \frac{\norm{g_k}^{4}}{2h_k}.
  \]
  By Cauchy--Schwarz and the Rayleigh-quotient bound,
  $h_k = g_k\trans B_k g_k \le \norm{B_k}\,\norm{g_k}^{2}$, hence
  \[
  f(x^{k}) - m_k(s^{k}_c)
    \ge \frac{\norm{g_k}^{2}}{2\norm{B_k}}
    \ge \tfrac12\norm{g_k}\,
            \min\left\{\Delta_k, \frac{\norm{g_k}}{\norm{B_k}}\right\}.
  \]
  
  \emph{Sub-case 2b: $\alpha_{\min} > \bar\alpha$.}\quad
  Then $\alpha^{k}_c = \bar\alpha = \Delta_k/\norm{g_k}$. Since $\psi$ is
  decreasing on $[0,\alpha_{\min})$, evaluating at $\bar\alpha$ and using
  $\bar\alpha < \alpha_{\min} = \norm{g_k}^{2}/h_k$, i.e.\
  $\bar\alpha\,h_k < \norm{g_k}^{2}$,
  \[
  f(x^{k}) - m_k(s^{k}_c)
    = \bar\alpha\norm{g_k}^{2} - \tfrac12\bar\alpha^{2}h_k
    > \bar\alpha\norm{g_k}^{2} - \tfrac12\bar\alpha\norm{g_k}^{2}
    = \tfrac12\bar\alpha\norm{g_k}^{2}
    = \tfrac12\Delta_k\norm{g_k}.
  \]
  This completes the proof.
\end{proof}

\begin{remark}
Lemma~\ref{lem:cauchy} is the key technical estimate underlying global
convergence: any step that does at least as well as the Cauchy step inherits
this lower bound on the model decrease, and combined with the ratio test in
GTR this drives $\norm{\grad f(x^{k})}\to 0$.
\end{remark}

\section{Global convergence of GTR}

\begin{theorem}[Global convergence of GTR]\label{thm:GTR-global}
Let $f\in\mathcal{C}^{2}(\R^{n})$ be bounded below with $\grad f$ $L$-Lipschitz continuous, and assume $\norm{B_{k}}$ is uniformly bounded. Apply GTR
(Algorithm~\ref{alg:GTR}) requiring $m_{k}(s^{k})\le m_{k}(s^{k}_{c})$. Then either
there exists $\ell\ge 0$ with $\grad f(x^{\ell})=0$, or
$\liminf_{k\to\infty}\norm{\grad f(x^{k})}=0$.
\end{theorem}
  
\medskip
The proof rests on two lemmas: the Cauchy-point decrease (Lemma~\ref{lem:cauchy})
and the lower bound on the trust-region radius below.
  
\begin{lemma}[Lower bound on the TR radius]\label{lem:tr-radius}
Under the assumptions of Theorem~\ref{thm:GTR-global}, suppose there exists
$\eps>0$ such that $\norm{\grad f(x^{k})}\ge\eps$ for all $k$. Then there
exists a constant $c\in(0,1)$ such that
\[
\Delta_{k}\;\ge\;\frac{c\,\eps}{L}\qquad\text{for all }k\ge 0.
\]
\end{lemma}

\begin{proof}
Write $g_{k}:=\grad f(x^{k})$. By the integral form of Taylor's theorem and $\norm{\grad^{2}f(x)}\le L$,
\[
  \bigl|f(x^{k}+s^{k})-m_{k}(s^{k})\bigr| \le \tfrac{L}{2}\norm{s^{k}}^{2} \le \tfrac{L}{2}\Delta_{k}^{2}.
\]
Combined with the Cauchy-point bound (Lemma~\ref{lem:cauchy}) and $\norm{B_{k}}=\norm{\grad^{2}f(x^{k})}\le L$,
\[
  f(x^{k})-m_{k}(s^{k}) \ge \tfrac{1}{2}\eps\, \min\left\{\Delta_{k},\,\tfrac{\eps}{L}\right\}.
\]
Hence
\[
  \bigl|\rho_{k}-1\bigr| = \frac{|f(x^k + s^k) - m_k(s^k)|}{f(x^k) - m_k(s^k)} \le \frac{(L/2)\Delta_{k}^{2}}{(\eps/2) \min\{\Delta_{k},\eps/L\}}.
\]

Now choose $c := \min\left\{\tfrac{1}{2},\,\tfrac{\Delta_{0}L}{\eps}\right\}\in(0,1)$. Suppose, for contradiction, that some iteration $k$ has $\Delta_{k}\le 2c\eps/L\le\eps/L$. Then $\min\{\Delta_{k},\eps/L\}=\Delta_{k}$, and substituting,
\[
|\rho_{k}-1| \le \frac{(L/2)\Delta_{k}^{2}}{(\eps/2)\Delta_{k}} = \frac{L\Delta_{k}}{\eps} \le 2c \le 1 - 0.9 + c \le 0.1.
\]
Therefore $\rho_{k}\ge 0.9$, so iteration $k$ is \emph{very successful} and the GTR rule sets $\Delta_{k+1}=2\Delta_{k}\ge\Delta_{k}$.

Thus the radius can never drop below $c\eps/L$: any iteration with
$\Delta_{k}\le 2c\eps/L$ at least preserves the radius, so by induction on $k$,
\[
\Delta_{k}\;\ge\;\min\!\left\{\Delta_{0},\,\tfrac{c\eps}{L}\right\}
            \;\ge\;\tfrac{c\eps}{L}
\]
for all $k\ge 0$ (after possibly shrinking $c$ further so that
$\Delta_{0}\ge c\eps/L$). This proves the claim.
\end{proof}
  
\begin{proof}[Proof of Theorem~\ref{thm:GTR-global}]
If $\grad f(x^{k})=0$ for some $k$, the algorithm terminates and there is
nothing to prove; this also covers the case of finitely many successful
iterations, since on every unsuccessful step $\Delta_{k+1}=\tfrac12\Delta_{k}$,
so $\Delta_{k}\to 0$, contradicting Lemma~\ref{lem:tr-radius} unless the
gradient at the last successful iterate is zero.

Suppose instead, for contradiction, that there exists $\eps>0$ with
$\norm{\grad f(x^{k})}\ge\eps$ for all $k$, and that the set
$\mathcal{S}:=\{k:\rho_{k}\ge 0.1\}$ of successful iterations is infinite.
For every $k\in\mathcal{S}$ the GTR update gives $x^{k+1}=x^{k}+s^{k}$ and
\[
f(x^{k})-f(x^{k+1})
    \;\ge\;0.1\bigl(f(x^{k})-m_{k}(s^{k})\bigr)
    \;\ge\;\tfrac{0.1}{2}\norm{g_{k}}\,
        \min\!\left\{\Delta_{k},\,\tfrac{\norm{g_{k}}}{\norm{B_{k}}}\right\},
\]
by Lemma~\ref{lem:cauchy}. Since $\norm{g_{k}}\ge\eps$ and
$\norm{B_{k}}=\norm{\grad^{2}f(x^{k})}\le L$,
\[
f(x^{k})-f(x^{k+1})
    \;\ge\;0.05\,\eps\,\min\!\left\{\Delta_{k},\,\tfrac{\eps}{L}\right\}
    \;\ge\;0.05\,\eps\,\min\!\left\{\tfrac{c\eps}{L},\,\tfrac{\eps}{L}\right\}
    \;=\;\tfrac{0.05\,c}{L}\,\eps^{2},
\tag{$\ast$}\label{eq:per-step}
\]
where the second inequality uses Lemma~\ref{lem:tr-radius} and $c\in(0,1)$.

Since $f(x^{k+1})=f(x^{k})$ on every unsuccessful iteration, telescoping the
decrease over all $k\le K-1$ gives
\[
f(x^{0})-f(x^{K})
    \;=\;\sum_{i=0}^{K-1}\!\bigl(f(x^{i})-f(x^{i+1})\bigr)
    \;=\;\sum_{i\in\mathcal{S},\,i<K}\!\bigl(f(x^{i})-f(x^{i+1})\bigr).
\]
Letting $K\to\infty$ and using $f(x^{K})\ge f_{\mathrm{low}}$,
\[
f(x^{0})-f_{\mathrm{low}}
  \;\ge\;\sum_{i\in\mathcal{S}}\bigl(f(x^{i})-f(x^{i+1})\bigr)
  \;\ge\;|\mathcal{S}|\cdot\tfrac{0.05\,c}{L}\,\eps^{2} \to \infty,
\tag{$\ast\ast$}\label{eq:telescoped}
\]
by~\eqref{eq:per-step}. But then left-hand side is finite by assumption, contradiction. Hence no such $\eps>0$ exists, and so $\liminf_{k\to\infty}\norm{\grad f(x^{k})}=0$.
\end{proof}

\begin{theorem}[Stronger global convergence]\label{thm:GTR-lim}
Under the conditions of Theorem~\ref{thm:GTR-global},
\[
  \lim_{k\to\infty}\norm{\grad f(x^{k})}=0.
\]
\qed
\end{theorem}

\section{Solving the trust-region subproblem}

The characterisation below shows that solving \eqref{eq:TRsub} reduces to the secular equation 
\[
  \norm{(B_{k}+\lambda I)^{-1}\grad f(x^{k})}=\Delta_{k}
\]
in the scalar $\lambda\ge\max\{0,-\lambda_{\min}(B_{k})\}$.

\begin{theorem}[Characterization of the global TR solution]\label{thm:TR-charac}
A vector $s^{*}\in\R^{n}$ is a global solution of \eqref{eq:TRsub} if and only if
$\norm{s^{*}}\le\Delta_{k}$ and there exists $\lambda^{*}\ge 0$ such that
\begin{enumerate}[leftmargin=2em,itemsep=0pt]
  \item $(B_{k}+\lambda^{*}I)s^{*}=-\grad f(x^{k})$,
  \item $\lambda^{*}(\Delta_{k}-\norm{s^{*}})=0$ \quad(complementarity),
  \item $B_{k}+\lambda^{*}I\succeq 0$.
\end{enumerate}
Moreover $\lambda^{*}$ is unique whenever $B_{k}+\lambda^{*}I\succ 0$. \qed
\end{theorem}

Theorem~\ref{thm:TR-charac} gives necessary and sufficient global
optimality conditions for a nonconvex optimization problem. It splits the practical solution into two cases:
\begin{itemize}[leftmargin=2em,itemsep=0pt]
  \item \textbf{Case 1:} $B_{k}\succ 0$ and the Newton step $s=-B_{k}^{-1}\grad f(x^{k})$ satisfies $\norm{s}\le\Delta_{k}$ — then $s^{*}=s$ and $\lambda^{*}=0$ (Newton step is inside the trust region).
  \item \textbf{Case 2:} otherwise — find $\lambda > \max\{0,-\lambda_{\min}(B_{k})\}$
        with $\norm{(B_{k}+\lambda I)^{-1}\grad f(x^{k})}=\Delta_{k}$, by a 1-D
        root-find; the solution $s^{*}$ lies on the boundary.
\end{itemize}

\begin{remark}[Large-scale problems]
For large $n$, repeated Cholesky factorisations of $B_{k}+\lambda I$ become
prohibitive. In practice one uses iterative methods (truncated conjugate gradient,
Lanczos / Steihaug--Toint) that produce an approximate solution at least as good
as the Cauchy point, hence preserving the global convergence guarantee.
\end{remark}

\begin{remark}[Linesearch vs.\ trust-region in practice]
Both linesearch and trust-region methods underpin state-of-the-art unconstrained
optimisation software. Linesearch needs more heuristics to deal with negative
curvature; TR handles indefinite Hessians naturally and accommodates quasi-Newton
$B_{k}$ that need not be positive definite. The choice between them is largely a
matter of taste.
\end{remark}


%==========================================================
\chapter{Optimality Conditions for Constrained Optimization}
%==========================================================

The unconstrained theory characterised optimality by $\grad f(x^{*})=0$. For
constrained problems this fails: the gradient at the minimiser need not vanish,
because feasibility is preventing further descent. The right generalisation is
the \emph{Karush--Kuhn--Tucker (KKT)} conditions, which characterise points at
which no \emph{feasible} descent direction exists.

The development proceeds in three steps:
\begin{enumerate}[leftmargin=2em,itemsep=0pt]
  \item Define geometric (tangent cone $T_{\Omega}$) and algebraic (linearised
        feasible directions $\mathcal{F}$) descriptions of feasible directions;
        always $T_{\Omega}\subseteq\mathcal{F}$, but the reverse can fail.
  \item Identify a \emph{constraint qualification (CQ)} ensuring
        $T_{\Omega}=\mathcal{F}$ — without a CQ, optimality conditions need not
        hold.
  \item Derive the KKT conditions from the algebraic description $\mathcal{F}$.
\end{enumerate}

\section{The constrained problem}

We consider the general constrained problem
\begin{equation}\label{eq:CP}\tag{CP}
\min_{x\in\R^{n}}\; f(x)\quad\text{subject to}\quad
c_{i}(x)=0,\;i\in E;\;\;c_{i}(x)\ge 0,\;i\in I,
\end{equation}
with feasible set
$\Omega=\{x\in\R^{n}:c_{i}(x)=0,\,i\in E;\;c_{i}(x)\ge 0,\,i\in I\}$,
and active set at $x\in\Omega$:
$\mathcal{A}(x)=E\cup\{i\in I:c_{i}(x)=0\}$.

\section{Tangent cone and linearized feasible directions}

\begin{definition}[Tangent cone]
The \emph{tangent cone} to $\Omega$ at $x\in\Omega$ is
\[
T_{\Omega}(x)=\Bigl\{s\in\R^{n}:\exists\,z^{k}\in\Omega,\,t^{k}\downarrow 0,\,
z^{k}\to x,\;s=\lim_{k\to\infty}\frac{z^{k}-x}{t^{k}}\Bigr\}.
\]
\end{definition}

\begin{definition}[Linearized feasible directions]
The set of \emph{linearized feasible directions} at $x\in\Omega$ is
\[
\mathcal{F}(x)=\{s:\grad c_{i}(x)\trans s=0,\,i\in E;\;
\grad c_{i}(x)\trans s\ge 0,\,i\in I\cap\mathcal{A}(x)\}.
\]
\end{definition}

\begin{lemma}\label{lem:TF}
$T_{\Omega}(x)\subseteq\mathcal{F}(x)$.
\end{lemma}
\begin{proof}\todo\end{proof}

\section{Constraint qualifications}

\begin{definition}[LICQ]
\eqref{eq:CP} satisfies the \emph{Linear Independence Constraint Qualification (LICQ)}
at $x\in\Omega$ if the active constraint gradients
$\{\grad c_{i}(x):i\in\mathcal{A}(x)\}$ are linearly independent.
\end{definition}

\begin{definition}[Slater CQ]
A convex problem \eqref{eq:CP} (with $c_{E}(x)=Ax-b$, $c_{i}$ concave for $i\in I$)
satisfies the \emph{Slater CQ} if there exists $\bar{x}$ with $A\bar{x}=b$ and
$c_{i}(\bar{x})>0$ for all $i\in I$.
\end{definition}

\begin{theorem}[CQ implies $T_{\Omega}=\mathcal{F}$]\label{thm:CQ}
Under LICQ (or any standard CQ), $T_{\Omega}(x)=\mathcal{F}(x)$.
\end{theorem}
\begin{proof}\todo\end{proof}

\begin{remark}[Linear constraints]
If the constraints of \eqref{eq:CP} are linear, i.e.\ $(c_{E},c_{I})(x)=Ax-b$,
then \emph{no constraint qualification is required}: the KKT conditions hold at
every local minimiser. This is the content of Theorem~17 in the slides.
\end{remark}

\section{First-order necessary conditions (KKT)}

\paragraph{Interpretation.} The KKT conditions generalise the unconstrained
stationarity condition $\grad f=0$: at a constrained minimiser, $\grad f$ lies in
the cone generated by the gradients of active constraints (with appropriate signs).
Inactive inequality constraints contribute nothing — by complementarity
$\lambda^{*}_{i}c_{i}(x^{*})=0$, so $c_{i}(x^{*})>0$ forces $\lambda^{*}_{i}=0$.

\begin{definition}[Lagrangian]
The \emph{Lagrangian} of \eqref{eq:CP} is
\[
\mathcal{L}(x,y,\lambda)=f(x)-\sum_{i\in E}y_{i}c_{i}(x)-\sum_{i\in I}\lambda_{i}c_{i}(x).
\]
\end{definition}

\begin{theorem}[KKT first-order necessary conditions]\label{thm:KKT}
Let $x^{*}$ be a local minimizer of \eqref{eq:CP} at which a CQ holds. Then
there exist multipliers $y^{*}\in\R^{|E|}$ and $\lambda^{*}\in\R^{|I|}$ such that
\begin{align*}
\grad f(x^{*})&=\sum_{i\in E}y^{*}_{i}\grad c_{i}(x^{*})+\sum_{i\in I}\lambda^{*}_{i}\grad c_{i}(x^{*}),\\
c_{i}(x^{*})&=0,\quad i\in E,\\
c_{i}(x^{*})&\ge 0,\quad i\in I,\\
\lambda^{*}_{i}&\ge 0,\quad i\in I,\\
\lambda^{*}_{i}c_{i}(x^{*})&=0,\quad i\in I \quad\text{(complementarity).}
\end{align*}
\end{theorem}
\begin{proof}\todo\end{proof}

\section{Second-order conditions}

Let $\mathcal{F}(\lambda^{*})=\{s\in\mathcal{F}(x^{*}):\grad c_{i}(x^{*})\trans s=0
\text{ for } i\in I\text{ with }\lambda^{*}_{i}>0\}$.

\begin{theorem}[Second-order necessary conditions]\label{thm:CP-SO-nec}
Let $x^{*}$ be a local minimizer of \eqref{eq:CP} at which a CQ holds, with
multipliers $(y^{*},\lambda^{*})$. Then
$s\trans\hess_{xx}\mathcal{L}(x^{*},y^{*},\lambda^{*})s\ge 0$ for all
$s\in\mathcal{F}(\lambda^{*})$.
\end{theorem}
\begin{proof}\todo\end{proof}

\begin{theorem}[Second-order sufficient conditions]\label{thm:CP-SO-suf}
Let $x^{*}$ be feasible for \eqref{eq:CP}, and let $(y^{*},\lambda^{*})$ satisfy
the KKT conditions at $x^{*}$. If
$s\trans\hess_{xx}\mathcal{L}(x^{*},y^{*},\lambda^{*})s>0$
for all $0\ne s\in\mathcal{F}(\lambda^{*})$, then $x^{*}$ is a strict local
minimizer of \eqref{eq:CP}.
\end{theorem}
\begin{proof}\todo\end{proof}

\section{Convex problems}

\begin{theorem}[KKT is sufficient for convex problems]\label{thm:CP-convex}
Suppose \eqref{eq:CP} is a convex problem (i.e.\ $f$ convex, $c_{i}$ concave for
$i\in I$, $c_{E}$ affine). If $\hat{x}$ is a KKT point of \eqref{eq:CP}, then
$\hat{x}$ is a global minimizer of \eqref{eq:CP}.
\end{theorem}
\begin{proof}\todo\end{proof}


%==========================================================
\chapter{Penalty Methods}
%==========================================================

Constrained optimisation introduces a fundamental tension: we want to minimise $f$
\emph{and} satisfy the constraints. For inequality-constrained or linear
equality-constrained problems we can sometimes maintain feasibility throughout,
but with general nonlinear equality constraints this is very hard, even
impossible. The idea behind \emph{penalty methods} is to replace the constrained
problem by a sequence of unconstrained problems whose minimisers approach the
original solution as a parameter is varied.

We now turn to algorithms for the equality-constrained problem
\begin{equation}\label{eq:eCP}\tag{eCP}
\min_{x\in\R^{n}} f(x)\quad\text{subject to}\quad c(x)=0,
\end{equation}
where $c=(c_{1},\dots,c_{m}):\R^{n}\to\R^{m}$ is smooth.

\section{The quadratic penalty function}

\begin{definition}[Quadratic penalty function]
For $\sigma>0$, the \emph{quadratic penalty function} associated with \eqref{eq:eCP} is
\begin{equation}\label{eq:Phi}
\Phi_{\sigma}(x)=f(x)+\frac{1}{2\sigma}\norm{c(x)}^{2}.
\end{equation}
\end{definition}

\paragraph{Interpretation.} The parameter $\sigma>0$ controls the penalty on
infeasibility. As $\sigma\to 0$, infeasibility is forced down and the unconstrained
minimisers of $\Phi_{\sigma}$ approach a solution of \eqref{eq:eCP}. The function
$\Phi_{\sigma}$ may have stationary points that are \emph{not} solutions of
\eqref{eq:eCP} — for example, local minimisers of $\norm{c(x)}$ when $c(x)=0$ has
no solution.

\begin{algorithm}[Basic quadratic penalty method]\label{alg:penalty}
Choose $\sigma_{0}>0$ and $\eps_{0}>0$. For $k=0,1,2,\dots$:
\begin{enumerate}[leftmargin=2em,itemsep=0pt]
  \item Find $x^{k}$ with $\norm{\grad\Phi_{\sigma_{k}}(x^{k})}\le\eps_{k}$.
  \item Choose $\sigma_{k+1}<\sigma_{k}$ and $\eps_{k+1}<\eps_{k}$.
\end{enumerate}
\end{algorithm}

\begin{lemma}[Gradient of $\Phi_{\sigma}$]\label{lem:gradPhi}
With $J(x)$ the Jacobian of $c$,
\[
\grad\Phi_{\sigma}(x)=\grad f(x)+\frac{1}{\sigma}J(x)\trans c(x)
=\grad f(x)-J(x)\trans y(x),
\]
where $y(x)=-c(x)/\sigma$.
\end{lemma}
\begin{proof}\todo\end{proof}

\section{Convergence of the penalty method}

\begin{theorem}[Global convergence of the quadratic penalty method]\label{thm:penalty}
Apply Algorithm~\ref{alg:penalty} to \eqref{eq:eCP}. Assume $f,c\in\mathcal{C}^{1}$, set
$y^{k}_{i}=-c_{i}(x^{k})/\sigma_{k}$, and assume
$\norm{\grad\Phi_{\sigma_{k}}(x^{k})}\le\eps_{k}$ with $\eps_{k}\to 0$ and
$\sigma_{k}\to 0$ as $k\to\infty$. Suppose further that $x^{k}\to x^{*}$ where
$\{\grad c_{i}(x^{*}):i=1,\dots,m\}$ are linearly independent. Then $x^{*}$ is a
KKT point of \eqref{eq:eCP} and $y^{k}\to y^{*}$, the corresponding Lagrange
multiplier vector.
\end{theorem}
\begin{proof}\todo\end{proof}

\begin{remark}
If $J(x^{*})$ is not full rank then $x^{*}$ instead (locally) minimizes the
infeasibility $\norm{c(x)}$.
\end{remark}

\begin{remark}[Ill-conditioning of $\hess\Phi_{\sigma_{k}}$]\label{rem:penalty-ill}
A serious practical drawback of the basic penalty method is that, asymptotically,
$\hess\Phi_{\sigma_{k}}(x^{k})$ has $m$ eigenvalues of order $O(1/\sigma_{k})$ and
the remaining $n-m$ eigenvalues of order $O(1)$ — so the condition number is
$O(1/\sigma_{k})$ and \emph{blows up as $\sigma_{k}\to 0$}. Consequently we may
not be able to compute $x^{k}$ accurately by Newton's method on $\Phi_{\sigma_{k}}$.
This motivates augmented Lagrangians (Chapter~\ref{ch:auglag}), which avoid the
need to drive $\sigma_{k}\to 0$.
\end{remark}

\section{Other penalty functions}

The quadratic penalty is \emph{smooth} but \emph{inexact}: $\sigma_{k}\to 0$ is
required, with the ill-conditioning consequences in
Remark~\ref{rem:penalty-ill}. \emph{Exact} penalty functions allow convergence at
a fixed (sufficiently small) $\sigma$, avoiding ill-conditioning, but are
typically non-smooth and so harder to minimise.

\begin{definition}[Exact penalty function]
A penalty function $\Phi(x,\sigma)$ is \emph{exact} if there exists $\sigma^{*}>0$
such that for all $\sigma<\sigma^{*}$, every local solution of the original problem
is a local minimizer of $\Phi(\cdot,\sigma)$.
\end{definition}

\begin{example}[$\ell_{1}$ and $\ell_{2}$ penalty functions]
For \eqref{eq:CP},
\[
\Phi_{1}(x,\sigma)=f(x)+\tfrac{1}{\sigma}\sum_{i\in E}|c_{i}(x)|+\tfrac{1}{\sigma}\sum_{i\in I}[c_{i}(x)]_{-},
\]
\[
\Phi_{2}(x,\sigma)=f(x)+\tfrac{1}{\sigma}\norm{c_{E}(x)}_{2}+\tfrac{1}{\sigma}\norm{[c_{I}(x)]_{-}}_{2}.
\]
Both are exact (but generally non-smooth); the quadratic penalty is smooth but
not exact.
\end{example}


%==========================================================
\chapter{Augmented Lagrangian Methods}
\label{ch:auglag}
%==========================================================

The quadratic penalty method has a fundamental drawback: $\sigma_{k}\to 0$ is
required, leading to severe ill-conditioning of the subproblem Hessian
(Remark~\ref{rem:penalty-ill}). The augmented Lagrangian method retains the
quadratic penalty term but adds an explicit estimate $u\in\R^{m}$ of the
Lagrange multipliers, so that as $u\to y^{*}$ feasibility is achieved
\emph{without} driving $\sigma$ to zero.

\section{The augmented Lagrangian}

\begin{definition}[Augmented Lagrangian]
The \emph{augmented Lagrangian} for \eqref{eq:eCP} is
\begin{equation}\label{eq:auglag}
\Phi(x,u,\sigma)=f(x)-u\trans c(x)+\frac{1}{2\sigma}\norm{c(x)}^{2},
\end{equation}
where $u\in\R^{m}$ is an estimate of the Lagrange multipliers and $\sigma>0$.
\end{definition}

\begin{remark}[Two interpretations of $\Phi(x,u,\sigma)$]
$\Phi(x,u,\sigma)$ admits two complementary readings:
\begin{itemize}[leftmargin=2em,itemsep=0pt]
  \item \textbf{Lagrangian augmented with a quadratic penalty}, $\mathcal{L}(x,u)+\frac{1}{2\sigma}\norm{c(x)}^{2}$.
        The quadratic term \emph{convexifies} the Lagrangian locally, fixing the
        well-known issue that the Lagrangian itself is generally a saddle in
        $(x,y)$ rather than a minimum in $x$.
  \item \textbf{Shifted quadratic penalty}: completing the square in $c(x)$ gives
        $\Phi(x,u,\sigma)=f(x)+\frac{1}{2\sigma}\norm{c(x)-\sigma u}^{2}-\tfrac{\sigma}{2}\norm{u}^{2}$,
        a quadratic penalty centred at the shifted target $c(x)=\sigma u$ rather
        than $c(x)=0$. As $u\to y^{*}$, the shift vanishes correctly even with
        $\sigma$ bounded away from zero.
\end{itemize}
The aim is to adjust both $u$ and $\sigma$ to encourage convergence.
\end{remark}

\paragraph{Lagrange multiplier estimates.} A short calculation gives
\[
\grad_{x}\Phi(x,u,\sigma)=\grad f(x)-J(x)\trans y(x), \qquad
y(x):=u-c(x)/\sigma,
\]
so the gradient of the augmented Lagrangian coincides with the gradient of the
ordinary Lagrangian evaluated at the multiplier estimate $y(x)$. This is the
update rule used in Algorithm~\ref{alg:auglag}.

\section{The basic algorithm}

\begin{algorithm}[Augmented Lagrangian method]\label{alg:auglag}
Choose $u^{0}\in\R^{m}$, $\sigma_{0}>0$, $\eps_{0}>0$. For $k=0,1,\dots$:
\begin{enumerate}[leftmargin=2em,itemsep=0pt]
  \item Find $x^{k}$ with $\norm{\grad_{x}\Phi(x^{k},u^{k},\sigma_{k})}\le\eps_{k}$.
  \item Update $u^{k+1}=u^{k}-c(x^{k})/\sigma_{k}$.
  \item Choose $\sigma_{k+1}\le\sigma_{k}$, $\eps_{k+1}<\eps_{k}$.
\end{enumerate}
\end{algorithm}

\section{Convergence}

\begin{theorem}[Global convergence of augmented Lagrangian]\label{thm:auglag}
Apply Algorithm~\ref{alg:auglag} to \eqref{eq:eCP}. Assume $f,c\in\mathcal{C}^{1}$,
let $y^{k}=u^{k}-c(x^{k})/\sigma_{k}$, and
$\norm{\grad_{x}\Phi(x^{k},u^{k},\sigma_{k})}\le\eps_{k}$ with $\eps_{k}\to 0$.
Suppose $x^{k}\to x^{*}$ where $\{\grad c_{i}(x^{*})\}_{i=1}^{m}$ are linearly
independent. Then $y^{k}\to y^{*}$ with
$\grad f(x^{*})-J(x^{*})\trans y^{*}=0$. Moreover, if either
$\sigma_{k}\to 0$ for bounded $u^{k}$, or $u^{k}\to y^{*}$ for bounded
$\sigma_{k}$, then $x^{*}$ is a KKT point of \eqref{eq:eCP} with associated
multipliers $y^{*}$.
\end{theorem}
\begin{proof}\todo\end{proof}

\begin{remark}[Why the augmented Lagrangian helps]
The crucial difference from the basic penalty method is the second alternative
in Theorem~\ref{thm:auglag}: convergence is achieved as soon as $u^{k}\to y^{*}$,
\emph{with $\sigma_{k}$ bounded away from zero}. This means the ill-conditioning
of the subproblem Hessian is controlled — $\hess_{xx}\Phi(x,u,\sigma)=
\hess_{xx}\mathcal{L}(x,y)+\frac{1}{\sigma}J(x)\trans J(x)$ has condition number
$O(1/\sigma)$, which stays bounded if $\sigma$ does. Augmented Lagrangian methods
underpin the LANCELOT software and are competitive with state-of-the-art
nonlinear programming codes.
\end{remark}


%==========================================================
\chapter{Interior Point Methods}
%==========================================================

For inequality-constrained problems, an alternative to penalty methods is to
\emph{stay strictly inside the feasible region throughout} and rely on a
\emph{barrier} that diverges as $c_{i}(x)\downarrow 0$. As we relax the barrier
parameter $\mu$, the unconstrained minimisers approach the boundary
along the so-called \emph{central path}.

We consider the inequality-constrained problem
\begin{equation}\label{eq:iCP}\tag{iCP}
\min_{x\in\R^{n}} f(x)\quad\text{subject to}\quad c(x)\ge 0,
\end{equation}
where $c=(c_{1},\dots,c_{p}):\R^{n}\to\R^{p}$ is smooth.

\section{The logarithmic barrier function}

\paragraph{Why the logarithm?} Since $-\log c_{i}(x)\to+\infty$ as $c_{i}(x)\downarrow 0$,
any minimiser $x(\mu)$ of $f_{\mu}$ must lie strictly inside $\Omega$ — the
constraint $c(x)>0$ is enforced \emph{implicitly} rather than as an explicit
inequality. When $\mu$ is large, $x(\mu)$ sits ``far inside'' $\Omega$; when $\mu$
is small, $f$ dominates and $x(\mu)$ is close to the optimal boundary of $\Omega$
(this is also the regime in which ill-conditioning kicks in; see below).

\begin{definition}[Barrier function]
For $\mu>0$, the \emph{logarithmic barrier function} is
\begin{equation}\label{eq:barrier}
f_{\mu}(x)=f(x)-\mu\sum_{i=1}^{p}\log c_{i}(x),
\end{equation}
defined on the strict interior $\{x:c(x)>0\}$.
\end{definition}

\begin{lemma}[Barrier derivatives]\label{lem:barrier-deriv}
With $J(x)$ the Jacobian of $c$ and $C(x)=\diag(c_{1}(x),\dots,c_{p}(x))$,
\begin{align*}
\grad f_{\mu}(x)&=\grad f(x)-\mu J(x)\trans c^{-1}(x)
=\grad f(x)-J(x)\trans \lambda(x),\\
\hess f_{\mu}(x)&=\hess f(x)-\sum_{i=1}^{p}\frac{\mu}{c_{i}(x)}\hess c_{i}(x)
+\mu J(x)\trans C(x)^{-2}J(x),
\end{align*}
where $\lambda_{i}(x):=\mu/c_{i}(x)$ are estimates of the Lagrange multipliers.
\end{lemma}
\begin{proof}\todo\end{proof}

\begin{remark}[Connection to the Lagrangian]
Lemma~\ref{lem:barrier-deriv} reveals a striking fact: $\grad f_{\mu}(x)=
\grad_{x}\mathcal{L}(x,\lambda(x))$, the gradient of the Lagrangian of
\eqref{eq:iCP} evaluated at the multiplier estimate $\lambda(x)=\mu/c(x)$.
Hence $f_{\mu}$ stationarity coincides with stationarity of the Lagrangian of
the original problem, with multipliers $\lambda(x)$. As $\mu\to 0$ the conditions
$\lambda_{i}(x)c_{i}(x)=\mu\to 0$ recover complementarity.
\end{remark}

\section{The central path}

\begin{theorem}[Local existence of the central path]\label{thm:central-path}
Assume $\Omega^{o}\ne\emptyset$ and $x^{*}$ is a local minimiser of \eqref{eq:iCP}
satisfying:
\begin{enumerate}[leftmargin=2em,itemsep=0pt]
  \item $\lambda^{*}_{i}>0$ whenever $c_{i}(x^{*})=0$ \emph{(strict
        complementarity)};
  \item $\{\grad c_{i}(x^{*}):i\in\mathcal{A}\}$ are linearly independent (LICQ);
  \item there exists $\alpha>0$ such that $s\trans\hess_{xx}\mathcal{L}(x^{*},\lambda^{*})s
        \ge\alpha\norm{s}^{2}$ for all $s$ with $J_{\mathcal{A}}(x^{*})s=0$.
\end{enumerate}
Then there exists a unique, continuously differentiable path
$\{x(\mu):\mu_{\eps}>\mu>0\}$ of (global) minimisers of $f_{\mu}$, with
$x(\mu)\to x^{*}$ as $\mu\to 0$.
\end{theorem}
\begin{proof}\todo\end{proof}

\section{The basic barrier algorithm}

\begin{algorithm}[Basic barrier algorithm]\label{alg:barrier}
Choose $\mu_{0}>0$, $\eps_{0}>0$ and $x^{0}$ with $c(x^{0})>0$.
For $k=0,1,\dots$:
\begin{enumerate}[leftmargin=2em,itemsep=0pt]
  \item Find $x^{k}$ with $c(x^{k})>0$ and $\norm{\grad f_{\mu_{k}}(x^{k})}\le\eps_{k}$.
  \item Choose $\mu_{k+1}<\mu_{k}$, $\eps_{k+1}<\eps_{k}$.
\end{enumerate}
\end{algorithm}

\section{Convergence}

\begin{theorem}[Global convergence of the barrier method]\label{thm:barrier}
Apply Algorithm~\ref{alg:barrier} to \eqref{eq:iCP}. Assume $f,c\in\mathcal{C}^{1}$,
$\lambda^{k}_{i}=\mu_{k}/c_{i}(x^{k})$ for $i=1,\dots,p$ with $c(x^{k})>0$,
$\mu_{k}>0$, and
$\norm{\grad f_{\mu_{k}}(x^{k})}\le\eps_{k}$ with $\eps_{k}\to 0$, $\mu_{k}\to 0$
as $k\to\infty$. Suppose $x^{k}\to x^{*}$ where the active gradients
$\{\grad c_{i}(x^{*}):i\in\mathcal{A}\}$ are linearly independent (i.e.\ LICQ),
$\mathcal{A}=\{i:c_{i}(x^{*})=0\}$. Then $x^{*}$ is a KKT point of
\eqref{eq:iCP} and $\lambda^{k}\to\lambda^{*}$, the associated Lagrange multipliers.
\end{theorem}
\begin{proof}\todo\end{proof}

\begin{remark}[Difficulties of basic (primal) IPMs]\label{rem:ipm-difficulties}\
\begin{itemize}[leftmargin=2em,itemsep=0pt]
  \item \emph{Ill-conditioning of $\hess f_{\mu}$.} As $\mu\to 0$,
        $\mu J(x)\trans C(x)^{-2}J(x)$ has $|\mathcal{A}|$ eigenvalues of order
        $O(1/\mu)$ tending to infinity, while the rest are $O(1)$. The condition
        number is $O(1/\mu_{k})$, blowing up. This is the main reason barrier
        methods fell out of favour in the 1960s, until primal-dual methods
        revived them in the 1980s.
  \item \emph{Poor starting points.} The current iterate $x^{k}$ is generally a
        poor starting point for the next subproblem $\min f_{\mu_{k+1}}$: a full
        Newton step typically violates $c(x^{k}+s^{k})>0$ as soon as $\mu$ is
        meaningfully decreased, so progress is slow.
\end{itemize}
The remedy for both is to switch to primal-dual methods.
\end{remark}

\section{Primal-dual interior point methods}

\begin{remark}[Primal-dual reformulation]
Modern interior point software (KNITRO, IPOPT, GALAHAD) works with the
\emph{perturbed KKT system}
\[
\grad f(x)-J(x)\trans\lambda=0,\qquad C(x)\lambda=\mu e,\qquad c(x)>0,\;\lambda>0,
\]
where $e=(1,\dots,1)\trans$. The path traced by solutions as $\mu\downarrow 0$ is
the \emph{central path}. Newton's method applied to the perturbed KKT system
treats $(x,\lambda)$ as joint variables — multiplier updates are computed
explicitly each iteration, in contrast to the implicit estimate
$\lambda(x)=\mu/c(x)$ in the primal method.

\textbf{Benefits over the primal barrier method:}
\begin{itemize}[leftmargin=2em,itemsep=0pt]
  \item Ill-conditioning is sidestepped: solving the linear system in the right
        $(x,\lambda)$ subspaces, the asymptotic conditioning is benign.
  \item Current iterates are good starting points for the next subproblem, so
        only a few Newton iterations per $\mu_{k}$ are needed.
  \item Aggressive parameter schedule $\mu_{k+1}=O(\mu_{k}^{2})$ yields
        \emph{superlinear} asymptotic convergence.
\end{itemize}
Computing an initial $x^{0}$ with $c(x^{0})>0$ is itself nontrivial; various
heuristics exist (e.g.\ phase-I problems).
\end{remark}

\begin{remark}[Linear programming: simplex vs.\ IPMs]
For LP, the simplex method has worst-case exponential complexity
(Klee--Minty 1972) but very good average-case behaviour. Karmarkar's IPM (1984)
gave the first practical polynomial-time algorithm for LP; the central path is
unique and global, and Newton's method on the barrier function can be precisely
quantified. IPMs scale to very large LPs, with numerically observed average
complexity $O(\log(\text{LP dim}))$ iterations, each more expensive than a
simplex pivot.
\end{remark}

%==========================================================
\appendix
%==========================================================

\chapter{Useful background}

\section{Taylor's theorem}

\begin{theorem}[First-order Taylor]
Let $f\in\mathcal{C}^{1}$. For any $x,s\in\R^{n}$ there exists $\tilde{\alpha}\in(0,1)$
with $f(x+s)=f(x)+\grad f(x+\tilde{\alpha}s)\trans s$.
\end{theorem}

\begin{theorem}[Second-order Taylor]
Let $f\in\mathcal{C}^{2}$. For any $x,s\in\R^{n}$ there exists $\tilde{\alpha}\in(0,1)$
with
\[
f(x+s)=f(x)+\grad f(x)\trans s+\tfrac12 s\trans\hess f(x+\tilde{\alpha}s)s.
\]
\end{theorem}

\section{Lipschitz continuity}

\begin{definition}
A function $g:\R^{n}\to\R^{m}$ is \emph{Lipschitz continuous} on $\R^{n}$
with constant $L>0$ if $\norm{g(y)-g(x)}\le L\norm{y-x}$ for all $x,y\in\R^{n}$.
\end{definition}

\section{Convexity}

\begin{theorem}[First-order characterization]
Let $f\in\mathcal{C}^{1}$. Then $f$ is convex if and only if
$f(y)\ge f(x)+\grad f(x)\trans(y-x)$ for all $x,y$.
\end{theorem}

\begin{theorem}[Second-order characterization]
Let $f\in\mathcal{C}^{2}$. Then $f$ is convex if and only if
$\hess f(x)\succeq 0$ for all $x$.
\end{theorem}

\end{document}